\section{The common marked datum and its rank-three cusp}

% BEGIN E72G1
\paragraph{E72G1. Both point tests and the finite end on one marked datum.}
Fix the original datum $D$, its coefficient DVR $O$, the given
$\iota$ and a prolongation $B$ above $\ell$. Set
$\eta=D.\eta$, $\Psi=D.\Psi$, $\sigma=D.\mathrm{cc}$ and
$M=\ell N|\operatorname{disc}K|$. The eight cohomological inputs
to AFE15 are supplied by G71H3, equivalently by the cohomological
projection of the rank-general hypothesis in M72S1. Apply AFE15
once. Keep its $L,g_1,\Phi,C_0$, its link $k$, and its
$\mathrm{Ok}$, $\mathrm{Big}$ and seed data. AFE17 gives the full
$\mathrm{Marks}(\eta,\Psi,\sigma,g_1)$.

To use the modularity conclusion of M72S2 at this particular link,
check the ramification premise of L69D3. The fixed field of
$\ker\eta$ in $\Omega$ is the embedded $K$, by $D.\eta_{\rm ker}$.
It is an imaginary quadratic field. By AFE5 it is contained in
$\mathbb Q(\mu_{|\operatorname{disc}K|})$. At any valuation for
which $M$ is a unit, the discriminant is a unit too. Reduction on
these roots of unity is injective, since their order is invertible
in the residue field. Inertia fixes their reductions and hence
fixes $K$. This proves that every valuation with inertia outside
$\ker\eta$ has $v(M)>0$ in additive notation. Also $\ell N\mid M$,
$e\ge 1$ and $N$ is even by the original datum. AFE1 supplies the
residue characteristic; $D.\mathrm{conj}$ supplies the specified
complex conjugation.

Apply L69D3 with $\mathrm{ModUAlt}$ from M72S2. On this very link
one has $\mathrm{FinPt}(O,\Phi,3)$,
$\mathrm{AutPtAlt}(k,3,\iota,\varepsilon_L^{-1})$, and a character
$a_2$ with $\mathrm{IsAlgChar}(\iota,a_2)$ and
$\mathrm{AutPtAlt}(k,2,\iota,a_2)$. The auxiliary extension of
coefficients in the two runs has already been removed from the
test-ring quantifiers in L69D3. AFE21 now gives
$\mathrm{FinEnd}(\ell,\eta,\Psi,\sigma,g_1,B)$ on the original
coefficients: its witness is the embedded $G_L$, chosen before
the test ring, its ordinary point and the element to test. The
marker and the two automorphic point tests have therefore been
fixed before any such specialization.
% END E72G1

% BEGIN E72G2
\paragraph{E72G2. Descent of the rank-three specialization on an arbitrary datum.}
Take the data of E72G1. Fix any allowed finite DVR subalgebra
$A\subset\overline{\mathbb Q}_\ell$ over $O$, any ordinary point
$r:G_{\mathbb Q}\to\mathrm{GG}_3(A)$, and the irreducibility
hypothesis on every open subgroup of $\ker\eta$ as specified in
$\mathrm{AutEnd3}$. Write $R(g)$ for the matrix component after
the given inclusion $A\hookrightarrow\overline{\mathbb Q}_\ell$.
Now choose any of the $K$-algebra structures on $\Omega$ allowed
by this test, and any $\rho$ with matrix $R$ on its $G_K$.
The quadratic extension $K/\mathbb Q$ is normal. For instance
each embedding sends a generator to one of two conjugate roots,
each of which generates the same subfield of $\Omega$. Hence
the image of this second embedding is the original $K$-subfield,
and its pointwise stabilizer is still $\ker\eta$.
The link identifies $G_L=G_{L^+}\cap\ker\eta$ inside
$G_{\mathbb Q}$. By Galois correspondence the image of $K$ lies
in the specified image of $L$ in $\Omega$. Taking the inverse
of the latter embedding on that subfield gives a $K\to L$
compatible with both selected embeddings. Since $L/\mathbb Q$
is abelian Galois by the link, $L/K$ in this actual tower is
finite abelian Galois.

Finiteness of $A/O$ implies integrality. Its maximal ideal lies
over that of the local ring $O$, so the coefficient inclusion is
local. AFE19 applied to $r$ gives the lift on $\Phi$ with the
matrix identity at every $s\in G_L$ and the three required inertia
characters with weights $(1,0,-1)$. The subgroup $G_L$ is open
in $G_{\mathbb Q}$ and lies in $\ker\eta$, so the stated matrix
irreducibility makes $\rho|_{G_L}$ irreducible. The test
$\mathrm{AutPtAlt}(k,3,\iota,\varepsilon_L^{-1})$ from E72G1
therefore yields a regular algebraic conjugate-self-dual cusp
attached to $(\rho|_{G_L})\otimes\varepsilon_L^{-1}$. Restriction
of the rational cyclotomic character through the compatible
tower identifies this with the restriction of
$t=\rho\otimes\varepsilon_K^{-1}$.

The descent uses a continuous representation with the required
polarization. The entries of $A$ are integral over $\mathbb Z_\ell$
and have norm at most one. Writing $\ell=u\varpi^b$ in the actual
test DVR, powers of its maximal ideal therefore map to arbitrarily
small norm balls. The stipulated adic continuity of $r$ thus gives
continuity of $\rho$, then of $t$ by continuity of the cyclotomic
character. For the pairing take the antidiagonal identity matrix $J$.
Normalization and the sign of $\sigma$ give $r(\sigma)=(1,-1)$.
For $s\in G_K$ the sign is $+1$ and the conjugation law in
$\mathrm{GG}_3$ reads
$R(\sigma s\sigma^{-1})=J(R(s)^{\rm t})^{-1}J$.
Writing $s^c$ for this conjugate in $G_K$ and using
$\varepsilon_K(s^c)=\varepsilon_K(s)$ gives
\[
 t(s)^{\rm t}Jt(s^c)
   =\varepsilon_K(s)^{-2}J
   =\varepsilon_K(s)^{\,1-3}J .
\]
The restriction of $t$ to $G_L$ stays irreducible since multiplication
by a scalar character preserves invariant subspaces.

Apply D70W1 successively along a cyclic prime-degree series for
$L/K$. Its hypotheses persist since the restriction at the very
top is irreducible. The result is a regular algebraic cusp over
the specified $K$ attached to this very $t$. By the Frobenius
calculation of D70W2 the multiplier $\varepsilon_K^{-2}$ gives the
dual Satake multisets at the conjugate primes, hence the cusp is
conjugate self-dual by strong multiplicity one. This is exactly
the original $\mathrm{Af.IsAutCsd}(\iota,t)$, whose definition
requires a cusp, regular algebraicity, conjugate duality and
cofinite attachment. Thus $\mathrm{AutEnd3}$ holds at the mark
of E72G1. The test permits the second embedding, every allowed
coefficient DVR, and every representation with the displayed
matrix identity.
% END E72G2
